\chapter{Notes} \label{chap:notes}


    \begin{definition}[Proof Analysis Problem]
        Given a proof system $Q$, the Proof Analysis Problem over $Q$ restricted to $s(n)$ is defined as:
        \[\mathrm{PAP}_Q[s(n)] = \{(\varphi, \pi, 1^t) \mid \varphi \in \mathrm{Sat}, \pi : Q \vdash \lnot \mathrm{Ref}^\mathrm{Res}(\varphi, s) \text{ with } t \geq s(n), \abs{\varphi} = n\}\]
    \end{definition}

    \begin{definition}[Proof Size Problem]
        Given a proof system $Q$, the Proof Size Problem over $Q$ is defined as:
        \[\mathrm{PSP}_Q = \{(\varphi, 1^t) \mid \exists \pi \text{ s.t. } \pi : Q \vdash \phi, \abs{\pi} \leq t\}\]
    \end{definition}

    For any proof system $Q$, the automatability of $Q$ can be viewed as some form of search version of $PSP_Q$, where $t$ is forced to be in function of $\abs{\varphi} + \mathrm{Size}_Q(\varphi)$.
    

    \begin{definition}[Automatability]
        A proof system $Q$ is said to be $t$-\textit{automatable} when there is an algorithm $\mathcal{A}$ such that $\mathcal{A}(\varphi)$ returns a $Q$-proof of the formula $\varphi$ in time at most $t(\abs{\varphi} + \mathrm{Size}_Q(\varphi))$.
    \end{definition}

    When $t = \mathrm{poly}$, we simply say that $Q$ is \textit{automatable}.


    \begin{definition}[Analyzability]
        A proof system $Q$ is said to be $(t,s)$-\textit{analyzable} when $\mathrm{PAP}_Q[s(n)] \in \mathsf{DTIME}(t(n))$.
    \end{definition}

    When $t,s = \mathrm{poly}(n)$, we simply say that $Q$ is \textit{analyzable}.


    It's immediate to see that automatability and analyzability can be combined to easily solve the satisfiability problem.

    \begin{proposition}
        Let $Q$ be a proof system such that $\mathrm{Res} \leq Q$. If $Q$ is $t$-automatable and $(t,t)$-analyzable then $\mathsf{NP} \subseteq \mathsf{DTIME}(t(n))$.
    \end{proposition}

    \begin{proof}
        Let $\mathcal{A}$ and $\mathcal{B}$ be two algorithms certifying, respectfully, that $Q$ is $t$-automatable and $(t,t)$-analyzable. Then, the algorithm $\mathcal{C}(\phi) = \mathrm{B}(\phi, \mathcal{A}(\phi), \abs{\phi}^2)$
    \end{proof}

    Under the standard complexity assumption that $\mathrm{NP}$ is not contained in $\mathsf{P}$ (or $\mathsf{QP}$, $\mathsf{SUBEXP}$), the properties of automatability and analyzability are disjoint when restricted to polynomial time algorithms (or quasi-polynomial, sub-exponential). We observe that, even if the above statement holds for proof systems that simulate Resolution, the disjointness of both properties under polynomial time algorithms seems to hold even for proof systems weaker than Resolution. For instance, de Rezende [dR20] proved that, under the Randomized Exponential Time Hypothesis (rETH), tree-like Resolution doesn't admit automating algorithms with runtime lower than quasi-polynomial.

\cleardoublepage